\documentclass[11pt,a4paper]{article}
\usepackage[T1]{fontenc}
\usepackage[utf8]{inputenc}
\usepackage{lmodern}
\usepackage{amsmath,amssymb}
\usepackage[hidelinks]{hyperref}
\title{Lightweight Temporal Anomaly Detection for Smart-Home Activities of Daily Living}
\author{}
\date{}
\begin{document}
\maketitle
\begin{abstract}
Do lightweight temporal models provide useful anomaly detection under missing sensors and household shift? This paper proposes a methodological evaluation, not a validated effect. We define a primary outcome, an independent validation strategy and explicit stopping criteria. Results are unavailable; theoretical assumptions and data eligibility must be audited before study activation.
\end{abstract}
\section{Research problem}
Do lightweight temporal models provide useful anomaly detection under missing sensors and household shift? The research is framed around a verifiable computational property rather than a change of application context alone. Its technical originality remains a hypothesis requiring a targeted literature review.
\section{Hypotheses and anticipated contribution}
H1 proposes that the candidate approach will improve the pre-specified primary outcome relative to a matched, non-trivial baseline. H2 proposes that the improvement, if observed, remains reliable under independent domains, fault conditions or held-out units. Neither hypothesis has been tested.
\section{Data, materials and eligibility}
Chronological and home-disjoint splits; evaluate label contamination and false alert burden. Data licensing, annotations, task comparability and reproducible split definitions are preconditions. Where authentic data are unavailable, only clearly labelled synthetic experiments are appropriate; simulated outcomes must never be reported as human or clinical measurements.
\section{Methods}
Benchmark rule-based, HMM and compact temporal neural models on licensed activity sequences with leave-home-out evaluation. All model selection and calibration must use training-only data; evaluation must avoid overlap between participants, sessions, households and temporal blocks as appropriate. Fixed seeds, environment manifests and source identifiers are required for repeatability.
\section{Evaluation protocol and statistics}
The primary outcome is: Event-level anomaly recall at a pre-specified false-alarm rate.. Analyses must specify the estimand, comparator, independent unit, confidence interval procedure, failure criteria and sensitivity checks before running confirmatory evaluation. Report both average performance and challenging edge cases.
\section{Results}
No empirical or simulated benchmark has yet been executed for this candidate publication. This section intentionally contains no numerical result, effect size, model comparison or inferred statistical significance.
\section{Discussion and validity limits}
SOURCE\_PLANNED\_ARTICLE: require verified household labels and a meaningful baseline; avoid clinical claims. A separate manuscript requires a distinct estimand and independently defendable novelty, not just a renamed application, sensor set or population.
\section{Conclusions}
This is a pre-experimental concept manuscript. Research activation requires an evidence-backed novelty decision, a frozen study protocol and an approved data or simulation pathway.
\end{document}
